\reportpage{\datacontinued}{
\subheading{四、同一次TCP上传连接（题11—20）}
以下均引用\tcpfile，\code{tcp.stream == 0}，使用相对Seq/Ack。C表示客户端，S表示服务器；TCP Len只计数据，不计TCP首部。
\ev{fig:syn}{syn-detail}{客户端SYN的序号、窗口和选项（帧1）}
\qa{11. 双方IP、端口及连接发起方}{
C为\code{10.135.133.233:65158}，S为\code{128.119.245.12:80}。帧1由C发送SYN、ACK=0，故C主动建立连接。端口见\figref{fig:syn}，对应的是\figref{fig:success}的实际文件上传。
}
\qa{12. 三次握手的Seq、Ack与标志}{
依据\figref{fig:syn}及最终文件帧1—3，记录如下：
\par\noindent\begin{tabularx}{\linewidth}{|c|c|c|c|c|c|X|}\hline
帧&方向&Seq&Ack&SYN&ACK&附加标志\\\hline
1&C$\to$S&0&0&1&0&ECE、CWR\\\hline
2&S$\to$C&0&1&1&1&ECE\\\hline
3&C$\to$S&1&1&0&1&无\\\hline
\end{tabularx}\par
SYN占用一个序号，因此两端确认SYN后Ack均为1。客户端与服务器原始ISN分别为3409073128、1221079522；帧3原始Seq/Ack为3409073129/1221079523。帧1的Ack字段虽为0，但ACK位未置位，其数值不作为有效确认。
}
}

\reportpage{\datacontinued}{
\ev{fig:synack}{synack-detail}{服务器SYN/ACK的窗口、MSS和缩放选项（帧2）}
\ev{fig:ackwindow}{ack-window}{第三次握手的相对序号和实际窗口（帧3）}
\qa{13. 双方通告的MSS}{
\figref{fig:syn}中客户端通告MSS 1460字节；\figref{fig:synack}中服务器通告1300字节。MSS表示本端希望接收的最大TCP数据段长度，约束对端发送，不包含TCP/IP首部。本次两个方向的通告值不同，不能把客户端的1460用作上传方向的接收上限。
}
\qa{14. 三次握手的Window与Window Scale}{
\par\noindent\begin{tabularx}{\linewidth}{|c|r|c|>{\raggedleft\arraybackslash}X|}\hline
帧&原始Window&本端通告shift&计算后窗口\\\hline
1&65535&6（$2^6=64$）&65535\\\hline
2&65160&7（$2^7=128$）&65160\\\hline
3&2053&无新选项&131392\\\hline
\end{tabularx}\par
帧1、2为带SYN的报文，其自身Window不缩放；不能直接将65535乘64或65160乘128作为该SYN报文的窗口。缩放已协商后，帧3客户端窗口为$2053\times64=131392$，与\figref{fig:ackwindow}相符。服务器后续帧14的原始Window为549，计算后$549\times128=70272$，见\figref{fig:rtt}。shift是左移位数，6和7不是直接乘数。
}
}

\reportpage{\datacontinued}{
\ev{fig:follow}{follow-stream}{Wireshark Follow TCP Stream：POST、文件名与正文}
\qa{15. 跟踪流中的应用层信息及上传验证}{
\figref{fig:follow}实际重组出的请求方法为\code{POST}，路径为\domain{/wireshark-labs/lab3-1-reply.htm}，Host为\domain{gaia.cs.umass.edu}。Content-Type为\code{multipart/form-data}；Content-Disposition中的\code{filename="alice.txt"}与所选公开文件一致，文件正文可见“ALICE'S ADVENTURES IN WONDERLAND”、作者Lewis Carroll以及开头“Alice was beginning to get very tired…”的内容。

HTTP Content-Length为152321，这是整个multipart请求体长度，不能误作文件长度。HTTP请求头678字节，总请求为$678+152321=152999$字节；实际文件为152138字节，HTTP首部与multipart封装额外占861字节。文件占客户端相对序号816—152953（闭区间）。

对最终流去重重组、剥离multipart封装后，文件与下载的\code{alice.txt}逐字节一致；SHA256均为\par\noindent\domain{d67c718cc4c14a864a5ffd9bc5555a2c0d74c54071f085e27f53b1bf84bc3d94}。

响应帧187为HTTP/1.1 200 OK；响应TCP数据780字节，其中HTML正文418字节，含Congratulations。该响应与\figref{fig:success}的浏览器页面相互印证；并非仅以页面点击动作判定上传成功。
}
}

\reportpage{\datacontinued}{
\ev{fig:data}{data-detail}{上传数据帧7的长度、序号及TCP选项}
\qa{16. 握手与数据段的TCP首部长度及选项差异}{
\figref{fig:syn}的客户端SYN首部44字节，即20字节固定首部加24字节选项；帧2为40字节。SYN中协商MSS、Window Scale、SACK Permitted，并含时间戳及NOP/EOL填充。\figref{fig:data}的普通数据段首部32字节，即20+12；选项为NOP、NOP、Timestamp，不再重复协商MSS和窗口缩放。首部长度由Data Offset以4字节为单位给出：11、10、8分别对应44、40、32。
}
\qa{17. 连续三个上传数据段的Seq和TCP Len}{
最终文件帧5、6、7连续发送，均为C$\to$S，Ack=1，记录如下，帧7界面见\figref{fig:data}。
\par\noindent\begin{tabularx}{\linewidth}{|c|r|r|r|X|}\hline
帧&Seq&TCP Len&下一Seq&相对时间/秒\\\hline
5&679&1288&1967&0.269128\\\hline
6&1967&1288&3255&0.269129\\\hline
7&3255&1288&4543&0.269131\\\hline
\end{tabularx}\par
满足下一Seq=当前Seq+数据长度。帧5前137字节为multipart前缀，余1151字节已是文件数据；帧6、7继续传送文件。实际1288小于服务器MSS 1300，且数据段含12字节TCP选项，不能把MSS直接当作每包实际数据长度。
}
}

\reportpage{\datacontinued}{
\ev{fig:rtt}{rtt-detail}{帧14累计确认帧7并显示269.348毫秒RTT}
\qa{18. 选择数据段与ACK并解释确认范围}{
选择\figref{fig:data}的帧7（C$\to$S，Seq=3255，Len=1288）与\figref{fig:rtt}的帧14（S$\to$C，Ack=4543）。帧7占序号3255—4542，Ack=4543表示下一个期望序号，累计确认了客户端序号1—4542，也包括之前帧4—6的数据；不表示已经收到第4543字节。

Wireshark的\code{tcp.analysis.acks\_frame}指向7，方向、四元组、序号区间与时间先后吻合。此处选择的是数据段的确认，不用握手ACK替代。
}
\qa{19. ACK RTT字段及手工时间差核验}{
\figref{fig:rtt}中“ACK to the segment in frame: 7”，并显示“RTT to ACK the segment: 269.348000 milliseconds”。对应字段\code{tcp.analysis.ack\_rtt}=0.269348秒。

帧7相对时间0.269131秒，帧14为0.538479秒，故\par
\noindent$0.538479-0.269131=0.269348\ \mathrm{s}=269.348\ \mathrm{ms}$。

这个样本含网络往返及接收方确认处理，可能受累计/延迟确认影响，不等于纯链路传播时延。它与SYN至SYN/ACK的268.199毫秒不是同一个样本；界面另列的iRTT也不能代替本题数据ACK的RTT。
}
}

\reportpage{\datacontinued}{
\ev{fig:release}{release-list}{帧184—190：Wireshark自定义列显示序号与关闭标志}
\qa{20. 谁先关闭、释放报文及实际合并情况}{
\figref{fig:release}显示服务器先在帧185发送FIN。下表按最终抓包的时间顺序记录，所有报文ACK位均为1，SYN为0；只有帧185、189的FIN为1。
\par\noindent\begin{tabularx}{\linewidth}{|c|c|r|r|r|X|}\hline
帧&方向&相对时间/s&Seq&Ack&内容\\\hline
184&S$\to$C&2.846628&781&153000&缺前段提示\\\hline
185&S$\to$C&6.876487&781&153000&FIN，Len 0\\\hline
186&C$\to$S&6.876715&153000&1&SACK，Len 0\\\hline
187&S$\to$C&7.197927&1&153000&HTTP 200，Len 780\\\hline
188&C$\to$S&7.198253&153000&782&ACK，Len 0\\\hline
189&C$\to$S&7.199853&153000&782&FIN，Len 0\\\hline
190&S$\to$C&7.466969&782&153001&ACK，Len 0\\\hline
\end{tabularx}\par
帧185的FIN占服务器序号781。帧186累计Ack仍为1，同时SACK区间为[781,782)，说明FIN所在区间已到，但服务器序号1—780尚未连续收到。帧187到达后，响应数据1—780与FIN781连成完整区间，帧188才累计Ack=782。

四个释放阶段可对应185（S FIN）、188（C ACK）、189（C FIN）、190（S ACK）。帧188、189实际为两包，没有合并。服务器FIN帧185同时以Ack=153000确认全部客户端请求数据；不能把“FIN报文也置ACK”误解成四次释放必然只有三包。帧190的Ack=153001确认了客户端FIN所占的一个序号。

\subheading{非理想现象与证据范围}
帧184提示Previous segment not captured，帧187标记TCP Retransmission；因此只能说响应数据晚于FIN被本机捕获，不能断言服务器先发送FIN再发送HTTP响应。该流有7帧重传/快重传、15帧重复ACK，无RST；Wireshark未标记\code{tcp.analysis.out\_of\_order}。采集日志丢弃计数为0，也不等于互联网中从未丢包。缺口经后续响应数据补齐，文件重组一致且双方正常关闭。
}
}
